% =====================================================================
% Introduction and problem statement  (author: Abhishek Singh, Stage 1)
% Prepared for the ELCVIA project report.
% =====================================================================
\section{Introduction}
\label{sec:intro}

Pneumonia is the single largest infectious cause of death in children worldwide; in 2019 it accounted for 14\% of all deaths of children under five~\cite{who_pneumonia}.
Chest X-rays are the most common imaging test for diagnosing it, but reading them takes trained radiologists,
who are scarce in many of the places where childhood pneumonia is most common.
Deep convolutional neural networks (CNNs) are therefore an attractive way to support diagnosis:
CheXNet~\cite{rajpurkar2017chexnet}, a 121-layer DenseNet, exceeded the average F1 score of four practising radiologists at detecting pneumonia
on frontal chest X-rays, and Kermany \etal~\cite{kermany2018identifying} showed that transfer learning can classify paediatric chest X-rays as
normal or pneumonia with high accuracy.

High accuracy alone does not show that a model is trustworthy. A CNN can reach a good score by learning \emph{shortcuts}~\cite{geirhos2020shortcut}:
features that happen to be correlated with the label in the training data but have nothing to do with the disease.
On chest X-rays, such shortcuts include text markers, laterality letters, image borders and hospital-specific acquisition details.
Zech \etal~\cite{zech2018variable} found that pneumonia models picked up hospital-specific information and generalised worse to other hospitals,
and DeGrave \etal~\cite{degrave2021ai} showed that COVID-19 detectors relied on markers and regions outside the lungs instead of lung findings.
A model that relies on such cues may fail without warning when it is used on images from a different scanner or hospital.

\paragraph{Problem statement.}
In this project we train a pneumonia classifier on paediatric chest X-rays and test \emph{where it looks}.
We ask two questions:
(i)~what fraction of a model's Grad-CAM~\cite{selvaraju2017gradcam} attention falls inside the lungs, and
(ii)~how removing everything outside the lungs (lung-masked training) changes accuracy and attention.

\paragraph{Approach and contributions.}
The work is organised in four stages that build on each other:
\begin{enumerate}
  \item A reproducible data pipeline with a \emph{patient-level} validation split, and a custom CNN baseline trained from scratch (\cref{sec:data}).
  \item DenseNet-121 transfer learning in the style of CheXNet, compared with the baseline.
  \item Grad-CAM heatmaps, lung masks from a pretrained segmentation model, and a lung-attention metric: the share of the heatmap inside the lung mask.
  \item Training on lung-masked images, a cross-test of the unmasked model on masked images, and a final comparison of accuracy and attention.
\end{enumerate}
  \paragraph{Main findings.}
  On the 624-image test set, the lung-masked DenseNet-121 achieved similar accuracy to the unmasked model (0.870 versus 0.873)
  and higher mean lung attention (51.79\% versus 39.52\%). When the unmasked model was evaluated on masked images, accuracy fell to 0.631.
  This indicates sensitivity to removing extra-pulmonary context, but does not identify a specific shortcut or establish causation.
Beyond the main experiments, our analysis of the raw data shows that image aspect ratio and colour mode differ between the two classes (\cref{sec:data-eda}).
Both are potential shortcuts that are present before any model is trained.
